% ****************************************************************************************** % Dissertation template and document class for Princeton University
% Author  : Jeffrey Scott Dwoskin <jdwoskin@princeton.edu>
% Adapted from: http://www.math.princeton.edu/graduate/tex/puthesis.html
% ****************************************************************************************** %


%%% For print copies
%% set 'singlespace' option to set entire thesis to single space, and define "\printmode" to remove all hyperlinks for printed copies of the thesis. Delete all output files before changing this mode -- it will turn hyperref package on and off
%\documentclass[12pt,lot, lof, singlespace]{puthesis}
%\newcommand{\printmode}{}

%%% For the electronic copy, use doublespacing, define "\proquestmode" to use outlined links, instead of colored links. 
\documentclass[12pt,lot, lof]{puthesis}
\newcommand{\proquestmode}{}
% I prefer proquestmode to be off for electronic copies for normal use, since the colored links are less distracting. However when printed in black and white, the colored links are difficult to read. 

%%% For early drafts without some of the frontmatter
% Also see the "ifodd" command below to disable more frontmatter
%\documentclass[12pt]{puthesis}

%%%%%%%%%%%%%%%%%%%%%%%%%%%%%%%%%%%%%%%%%%%%%%%%%%%%%%%%%%%%%\
%%%% Author & title page info

\title{An integrative `omics approach to quantitative analysis of physiological metabolic flux control}

\submitted{November 2015}  % degree conferral date (January, April, June, September, or November)
\copyrightyear{2015}  % year in which the copyright is secured by publication of the dissertation.
\author{Sean Richard Hackett}
\adviser{Professor Josh Rabinowitz}  %replace with the full name of your adviser
\departmentprefix{Program in}  % defaults to "Department of", but programs need to change this.
\department{Quantitative and Computational Biology}

%%%%%%%%%%%%%%%%%%%%%%%%%%%%%%%%%%%%%%%%%%%%%%%%%%%%%%%%%%%%%\
%%%% Tweak float placements
% From: http://mintaka.sdsu.edu/GF/bibliog/latex/floats.html "Controlling LaTeX Floats"
% and based on: http://www.tex.ac.uk/cgi-bin/texfaq2html?label=floats
% LaTeX defaults listed at: http://people.cs.uu.nl/piet/floats/node1.html

% Alter some LaTeX defaults for better treatment of figures:
    % See p.105 of "TeX Unbound" for suggested values.
    % See pp. 199-200 of Lamport's "LaTeX" book for details.
    %   General parameters, for ALL pages:
    \renewcommand{\topfraction}{0.85}	% max fraction of floats at top
    \renewcommand{\bottomfraction}{0.6}	% max fraction of floats at bottom
    %   Parameters for TEXT pages (not float pages):
    \setcounter{topnumber}{2}
    \setcounter{bottomnumber}{2}
    \setcounter{totalnumber}{4}     % 2 may work better
    \setcounter{dbltopnumber}{2}    % for 2-column pages
    \renewcommand{\dbltopfraction}{0.66}	% fit big float above 2-col. text
    \renewcommand{\textfraction}{0.15}	% allow minimal text w. figs
    %   Parameters for FLOAT pages (not text pages):
    \renewcommand{\floatpagefraction}{0.66}	% require fuller float pages
	% N.B.: floatpagefraction MUST be less than topfraction !!
    \renewcommand{\dblfloatpagefraction}{0.66}	% require fuller float pages

% The documentclass already sets parameters to make a high penalty for widows and orphans. 

%%%%%%%%%%%%%%%%%%%%%%%%%%%%%%%%%%%%%%%%%%%%%%%%%%%%%%%%%%%%%\
%%%% Use packages

%\usepackage{amsfonts}

%%% For figures
\usepackage{graphicx}
%\usepackage{subfig,rotate}

%%% for comments
\usepackage{verbatim}

%%% For tables
\usepackage{multirow}
% Longtable lets you have tables that span multiple pages.
\usepackage{longtable}

% Booktabs produces far nicer tables than the standard LaTeX tables.
%   see: http://en.wikibooks.org/wiki/LaTeX/Tables
\usepackage{booktabs}

%set parameters for longtable:
% default caption width is 4in for longtable, but wider for normal tables
\setlength{\LTcapwidth}{\textwidth}



%%%%%%%%%%%%%%%%%%%%%%%%%%%%%%%%%%%%%%%%%%%%%%%%%%%%%%%%%%
%%% Printed vs. online formatting
\ifdefined\printmode

% Printed copy
% url package understands urls (with proper line-breaks) without hyperlinking them
\usepackage{url}


\else

\ifdefined\proquestmode
%ProQuest copy -- http://www.princeton.edu/~mudd/thesis/Submissionguide.pdf

% ProQuest requires a double spaced version (set previously). They will take an electronic copy, so we want links in the pdf, but also copies may be printed or made into microfilm in black and white, so we want outlined links instead of colored links.
\usepackage{hyperref}
\hypersetup{bookmarksnumbered}

% copy the already-set title and author to use in the pdf properties
\makeatletter
\hypersetup{pdftitle=\@title,pdfauthor=\@author}
\makeatother

\else
% Online copy

% adds internal linked references, pdf bookmarks, etc

% turn all references and citations into hyperlinks:
%  -- not for printed copies
% -- automatically includes url package
% options:
%   colorlinks makes links by coloring the text instead of putting a rectangle around the text.
\usepackage{hyperref}
\hypersetup{colorlinks,bookmarksnumbered}

% copy the already-set title and author to use in the pdf properties
\makeatletter
\hypersetup{pdftitle=\@title,pdfauthor=\@author}
\makeatother

% make the page number rather than the text be the link for ToC entries
%\hypersetup{linktocpage}
\fi % proquest or online formatting
\fi % printed or online formatting


%%%%%%%%%%%%%%%%%%%%%%%%%%%%%%%%%%%%%%%%%%%%%%%%%%%%%%%%%%%%%\
%%%% Define commands

% Define any custom commands that you want to use.
% For example, highlight notes for future edits to the thesis
%\newcommand{\todo}[1]{\textbf{\emph{TODO:}#1}}


% create an environment that will indent text
% see: http://latex.computersci.org/Reference/ListEnvironments
% 	\raggedright makes them left aligned instead of justified
\newenvironment{indenttext}{
    \begin{list}{}{ \itemsep 0in \itemindent 0in
    \labelsep 0in \labelwidth 0in
    \listparindent 0in
    \topsep 0in \partopsep 0in \parskip 0in \parsep 0in
    \leftmargin 1em \rightmargin 0in
    \raggedright
    }
    \item
  }
  {\end{list}}

% another environment that's an indented list, with no spaces between items -- if we want multiple items/lines. Useful in tables. Use \item inside the environment.
% 	\raggedright makes them left aligned instead of justified
\newenvironment{indentlist}{
    \begin{list}{}{ \itemsep 0in \itemindent 0in
    \labelsep 0in \labelwidth 0in
    \listparindent 0in
    \topsep 0in \partopsep 0in \parskip 0in \parsep 0in
    \leftmargin 1em \rightmargin 0in
    \raggedright
    }

  }
  {\end{list}}



%%%%%%%%%%%%%%%%%%%%%%%%%%%%%%%%%%%%%%%%%%%%%%%%%%%%%%%%%%%%%\
%%%% Front-matter

% For early drafts, you may want to disable some of the frontmatter. Simply change this to "\ifodd 1" to do so.
\ifodd 0
% front-matter disabled while writing chapters
\renewcommand{\maketitlepage}{}
\renewcommand*{\makecopyrightpage}{}
\renewcommand*{\makeabstract}{}

% you can just skip the \acknowledgements and \dedication commands to leave out these sections.

\else


\abstract{
% Abstract can be any length, but should be max 350 words for a Dissertation for ProQuest's print indicies (150 words for a Master's Thesis) or it will be truncated for those uses.
Advances in genomics continue to make detection of large numbers of related biomolecules (e.g. mRNAs, proteins or metabolites) more routine. While these �omic techniques summarize how the state of one type of biomolecule changes across conditions, they provide a one-dimensional view of cellular processes that emerge at the interface between classes of species. To bridge multiple �omic dimensions, approaches are needed that rationally integrate diverse types of data. In assessing this question, we are interested in testing the validity of models that posit some interaction between diverse chemical species. This process entails not only a rigorous analysis of individual �omic datasets, but also a procedure for posing alternative models and statistically evaluating their relative support. Here, we present a quantitative, scalable strategy for revealing steady-state metabolic regulation by integrating metabolomic, proteomic, and flux measurements. This involves analyzing, on a reaction-by-reaction basis, whether fluxes across conditions are accounted for by a Michaelis-Menten relationship between enzymes, substrates, and potential regulator concentrations. We collected the required data for yeast growing at 25 different nutrient-limited steady states and applied this strategy to reveal the primary physiological flux control mechanisms for over 40 metabolic reactions, encompassing 34 instances of physiologically-relevant allosteric regulation. The identified regulation included classical feedbacks and unexpected cross-pathway connections. Quantitatively, half of flux control resided in substrates and one quarter in enzyme concentrations. For reversible reactions, the remainder resided largely in product levels, and for irreversible reactions, in allosteric effectors. Thus, metabolic activity is substantially self-regulated by metabolites themselves. Across the diverse growth conditions studied, strong changes in mRNA expression generally resulted in corresponding changes in protein abundance. This association between mRNA and protein expression changes, however, was far weaker than expected; the levels of many proteins departed markedly from their cognate transcript. These deviations were highly non-random, suggesting that post-transcriptional regulation has an important role in modulating the cellular response to nutrient availability. This work collectively provides focused examples of how the structures of complex systems can be interrogated to identify meaningful regulatory relationships.  And, it is at the interface between �omic datasets that these valuable relations emerge.
}

\acknowledgements{
%I would like to thank...
I am very grateful to the members of the Rabinowitz, Storey and Botstein groups as well as other members of the Lewis-Sigler Institute for providing contrasting perspectives that naturally support interdisciplinary thinking. Working in this environment has refined my scientific interests and set me on a lifelong path of probing inspiring scientific questions.

I would like to begin by thanking my advisor, Josh Rabinowitz, whose support and direction have shaped my project immensely. Josh paralleled my interest in metabolic systems biology, providing motivation and direction towards a meaningful, shared goal. A masterful communicator, Josh has refined the scope of my scientific expression: verbally, visually and linguistically.

John Storey can also not go without acklowledgement; John introduced me to statistical methods and approaches of visualization which revolutionized my approach to research. Being a fly on the wall at his group meetings helped me think about questions in a statistical manner. 

I would like to thank all past and present members of the Rabinowitz research group for their advice and support over the past five years. I would like to further highlight the formative contributions of several individuals. Vito Zanotelli provided enormous support by establishing a formalism for analyzing reaction kinetics. He helped assemble summaries of literature regulation and implemented approaches to deal with flux uncertainty. Wenxin Xu has helped with verifying regulatory predictions and testing the validity of flux estimation methods. Chel Nofal has been an invaluable sounding board for evaluating the merit and delivery of results. Ian Lewis assisted with the experimental measurement of media samples using NMR. I am also extremely grateful to Greg Ducker, Lukas Tanner, Xiaoyang Su, Yifan Xu, Chris Crutchfield and Meytal Higgens for helping with experimental methods.

Other members of the Lewis-Sigler Institute have enormously helped my work and development at Princeton: David Bostein set me on the path to thinking about biology in a quantitative manner and helped me to understand the awesome power of yeast genetics. Pat Gibney taught me how use yeast genetics; he made several of the expression strains used to test regulation and provided invaluable experimental advice on all things yeast. Dave Robinson provided programming advice and was always available and ready to think about statistical questions with me. Jonathan Goya helped run many of the chemostat experiments, was an integral part of the proteomics analysis and helped me think about the quantitative relationship between proteins and transcripts. David Perlman provided a plethora of advice on setting up proteomics experiments and spent a huge amount of effort to ensure that the dataset was top notch. Sandy Silverman taught me how to do numerous yeast experiments. I would also like to thank Keyur Desai and Peter Andolfatto for their guidance.

I would like to thank the great administrators whose helped me easily navigate bureaucratic challenges, including Marybeth Fidele, Jen Brick and Ping Ge.

I would also like to sincerely thank my past scientific mentors, who set me on the path to a PhD by making me see the intrigue of science. As an undergraduate, Teresa Gunn helped turn my interest in genetics into an interest in scientific discovery. Andy Clark taught me to think about grand questions and how to interrogate them using appropriate large datasets. Tony Greenberg fostered my interest in statistics, programming (and Brazillian Jiu-Jitsu).

I would also like to thank my family. As the son of two veterinarians,  I was instilled with a keen interest in biology that provided me direction through college and beyond. Finally, I want thank my fianc\'{e}, Maya, whose presence gave me perspective throughout my graduate education and whose proofreading efforts will hopefully make the reading of this dissertation more enjoyable.

This research is based upon work supported by the U.S. Department of Energy, Office of Science, Office of Biological and Environmental Research (BER), under Award Number DE-SC0012461. This research was supported in part by an award from the Department of Energy (DOE) Office of Science Graduate Fellowship Program (DOE SCGF). The DOE SCGF Program was made possible in part by the American Recovery and Reinvestment Act of 2009.  The DOE SCGF program is administered by the Oak Ridge Institute for Science and Education for the DOE. ORISE is managed by Oak Ridge Associated Universities (ORAU) under DOE contract number DE-AC05-06OR23100.  All opinions expressed in this work are the author's and do not necessarily reflect the policies and views of DOE, ORAU, or ORISE.
}

%\dedication{}

\fi  % disable frontmatter


%%%%%%%%%%%%%%%%%%%%%%%%%%%%%%%%%%%%%%%%%%%%%%%%%%%%%%%%%%%%%\
%%%% Hide some chapters

%%% If you want to produce a pdf that includes only certain chapters, specify them with includeonly, in addition to including all chapters below.
%\includeonly{ch-intro/chapter-intro}
%%% You can also specify multiple chapters.
%\includeonly{ch-intro/chapter-intro,ch-usage/chapter-usage}
%\includeonly{chap1,chap2,chap3}


%%%%%%%%%%%%%%%%%%%%%%%%%%%%%%%%%%%%%%%%%%%%%%%%%%%%%%%%%%%%%
%%%% Notes:

% Footnotes should be placed after punctuation.\footnote{place here.}
% Generally, place citations before the period~\cite{anotherauthor}.
% The proper usage for i.e., and e.g., include commas ``(e.g., option A, option B)''

%%%%%%%%%%%%%%%%%%%%%%%%%%%%%%%%%%%%%%%%%%%%%%%%%%%%%%%%%%%%%
%%%% Import chapters

\begin{document}

\makefrontmatter


% If you've disabled frontmatter, you can insert the toc manually
%\tableofcontents\clearpage

% \include lets us split up the document (and each include starts a new page):

\chapter{Introduction\label{ch:intro}}

As an undergraduate genetics major I became fascinated with complex traits and diseases such as height, corn yield and heart attack risk. While these traits are of huge societal importance, each is collectively determined by numerous genetic variants and environmental factors and consequently relating quantitative traits to their genetic determinants has been notoriously difficult \cite{Manolio:2009jp}. As I began at Princeton, I wanted to better understand the source of missing heritability and how the explanatory insufficiency of the GWAS paradigm may be better addressed. My insight into this problem was that the gulf separating genetic variation and whole-body complex traits is the massive complexity of organismal and cellular physiology.

At the scale of cellular and organism-level systems, thousands of genes and environmental effects collide leading to non-additive, non-linear interactions, such as those resulting in canalization \cite{Waddington:1942wy}. Because the outputs of these systems cannot be explained using additive genetics, the combination of many such sub-systems into organism-level phenotypes will only further distort the relationship between individual genetic variants and overall phenotypic variability. Outside of tightly controlled studies of microbial growth \cite{Bloom:2013bq} it is likely that additive genetics can only explain a minority of many phenotype's variation regardless of statistical power \cite{Weedon:2008gc,WellcomeTrustCaseControlConsortium:2007do}. Furthermore, genome-scale methods which attempt to account for such non-additive and non-linear interactions will be unable to rigorously detect such interactions because of the massive space of such interactions \cite{Friedman:1997kn}. An alternative to this unsupervised approach can be found in systems biology, where researchers continually build bottom-up models of systems operations based on a relatively small number of inputs.

\subsection{Harnessing the power of systems biology}

A single systems biology model may focus on a minute aspect of an organism, such as how the expression of a single gene, your favorite gene (YFG), is determined based on the activity of its primary transcriptional repressor and cis-regulatory variation in YFG \cite{Nuzhdin:2012ii}. Such a simple model of YFG expression can be used to characterize YFG protein abundance \cite{Jovanovic:2015hp}, which can in turn be built into a model of signaling or metabolic pathway activity \cite{Neves:2002bk, Chassagnole:2002ty}.  The importance of each model is that it integrates information within the appropriate context of the systems operation, allowing changes in system output (and uncertainty) to be expressed as a non-linear function of variation in inputs.

In a perfect world where systems are well understood, genetic and environmental variants could be propagated through individual subsystems, appropriately aggregating information and ultimately allowing for accurate prediction of quantitative traits and disease risk.  Importantly, at such a juncture, prediction would require far less data than would be necessary to initially parameterize the system. This utopian vision of systems biology is quite removed from reality, as the creation of systems-level models is hindered by both extensive experimental requirements and general challenges in model identification.

The need for extensive experimentation and challenges in model discovery lay at odds in experimental design, as given fixed resources, there are tradeoffs between producing an experiment which measures all of the desired variables and characterizing a sufficient number of conditions (and replicates) to distill anything meaningful from the quantitative information. This problem is further exacerbated as funding is further distributed among groups characterizing different systems, in different organisms, with limited direct cross-domain translatability. 

To simplify the complex model inference problems of systems biology, I adopted an experimentally-guided reductionist approach. Essentially, by acquiring more data, we can actually get away with analyzing fewer distinct conditions, by breaking complicated systems-level model inference questions into simpler sub-problems. To illustrate this reoccurring theme, consider the inference of metabolic regulation, the major topic of \hyperref[ch:simmer]{Chapter \ref{ch:simmer}}. If we are interested in identifying which of 50 metabolites ($m$) regulates each of 50 reactions of interest ($n$), then 2500, $\mathcal{O}(nm)$, possible metabolite-reaction pairs exist. This scale of inference is highly tractable, however, using existing approaches metabolic regulation is primarily tested at a systems-level (i.e. as a combined model-fit over all 50 reactions) \cite{Link:2013dj, Zampar:2013fr}.  In this case, the number of sets of possible regulatory relationships that could exist is 8.88 $\times 10^{84}$, $\mathcal{O}(n^{m})$, similar to the number of atoms in the observable universe. To make this problem scalable, I was able to achieve $\mathcal{O}(nm)$ complexity by breaking the dependence between reactions through measuring metabolic fluxes, rather than treating them as latent variables.

\subsection{Towards understanding how metabolism is controlled}

To allow for the application of piecewise inference, I focused on studying the baker's yeast, \textit{S. cerevisiae}, as a simple eukaryote that can be easily cultured and is amenable to genetic manipulation. Previous research in the Rabinowitz and Botstein groups has shown that when yeast are cultured in chemostats with greatly varying nutrient availability, both the transcriptome and metabolome are radically affected, but in a distinct manner \cite{Brauer:2008jn, Boer:2010fb}. Across these cultures, differences in growth rate and which nutrient is limiting growth (such as a carbon or nitrogen source), yields a panel of diverse physiological and metabolic states. As growth rates change, the transcriptome is strongly affected, with many transcripts increasing or decreasing in a log-linear fashion with growth rate, irrespective of which nutrient limits growth \cite{Brauer:2008jn}. In contrast, the metabolome is profoundly affected by which nutrient limits growth, with distinct metabolites appearing to be growth-limiting or overflowing depending on which primary nutrient limits growth \cite{Boer:2010fb}. While these studies convey the importance of both transcription and metabolism in nutrient-dependent growth control, each dataset is only able to provide a 1-dimensional answer to a question that inherently involves multiple `omic dimensions.

When yeast are growing at different rates and in media of greatly differing composition, they face a massive metabolic challenge. When growing quickly, metabolic rates (fluxes) must increase to produce enough  proteins, RNAs and other biomass components to keep pace with dilution.  Moreover, metabolic pathways must generally maintain an appropriate balance of energy, reducing equivalents and biomass components, across a broad spectrum of possible environments and genetic states. To achieve such fine control, pathway fluxes are collectively determined by enzyme expression, nutrient availability and metabolic regulation.  While analysis at this scale has been used to investigate the coordinated regulation of specific metabolic events involving small sets of genes \cite{Zampar:2013fr, Link:2013dj} as previously discussed this approach is ill suited to identifying novel regulation and reaction kinetics at scale. Instead, pathway-level regulation can be considered as an emergent property of interactions between reactions, where reaction-level kinetics are sufficient to reconstruct higher-order behavior \cite{Fell:1997wg}.

To characterize individual reactions, we can assess whether a model of reaction kinetics, such as Michaelis-Menten kinetics (\hyperref[intoEqtn:mm]{Equation \ref{intoEqtn:mm}}), is able to accurately predict the model output (flux) as a function of model inputs (metabolite and enzyme concentrations) and kinetic parameters \cite{Anonymous:1913wn, Liebermeister:2006fm, Tummler:2014cp}. Here, $\nu$ is the reaction flux, $\left[E\right]$ is the enzyme's concentration, k$_{cat}$ is the per-enzyme maximum rate of conversion of substrates into products ($\sfrac{1}{s}$), $\left[S\right]$ is the substrate's concentration, and k$_{m}$ is effectively the concentration where half of enzymes are bound by substrates.

\begin{equation}
\nu = k_{cat}\left[E\right]\frac{\left[S\right]}{\left[S\right] + k_{m}}\label{intoEqtn:mm}
\end{equation}

Under this simplified kinetic model, variable flux could be accomplished through two simple regulatory paradigms: either the concentration of enzymes changes proportionally to their flux, or the fractional substrate occupancy of each enzyme increases. For this second hypothesis, the upstream accumulation of metabolites could be a driving force that is directly controlled by nutrient availability. A combination of these strategies is also possible, where different enzymes could employ one or both of these methods of flux-control. 

Previous work from our laboratory, using transcript levels as an approximation of protein levels, has revealed that neither of these simple paradigms is generally consistent with inferred reaction fluxes \cite{Bradley:2009fj}. There is no consistent relationship between fluxes and transcript abundance for most enzymes, nor do metabolites reliably accumulate upstream of enzymes as flux increases. This suggests that either transcript abundance is an unsuitable surrogate for enzyme abundance (a hypothesis that can be tested using protemics) or that we need to expand the model beyond simple Michaelis-Menten kinetics to account for other mechanisms of flux regulation, such as post-translational modifications, allostery and localization. Across different nutritional conditions, we must consider not just the differences in flux associated with growth rate, but also potential differences in flux due variable cellular composition and energetics, such as tradeoffs between fermentation and respiration \cite{Lange:2001th, Feist:2010hq, BARFORD:1979ei}.

\subsection{Integrative `omic approaches to identify metabolic and post-transcriptional regulation}

While analysis of transcriptional and metabolomic variation suggested that existing yeast chemostat conditions would be an ideal system to investigate the regulatory underpinnings of nutrient-based growth control, additional data was required if we were to meaningfully address this question. To greatly reduce the gaps in our knowledge, I measured protein abundance across the 25 chemostat conditions using quantitative proteomics and estimated metabolism-wide fluxes using flux balance analysis (\hyperref[introFig:primarytopics]{Figure \ref{introFig:primarytopics}A}) \cite{Orth:2010hb}. As metabolite and enzyme concentrations are the primary determinants of reaction flux, the appropriate quantification of these species using mass spectrometry is a major focus of my research. During \hyperref[ch:quant_anal]{Chapter \ref{ch:quant_anal}} I discuss appropriate methods for analyzing mass spectrometry data and highlight my collaborations where metabolomics or proteomics provided novel insights into physiology and evolution. 

During \hyperref[ch:simmer]{Chapter \ref{ch:simmer}}, I discuss the integration of reaction flux with enzyme and metabolite concentrations to test alternative models of reaction-level kinetics. The first aim of this study is determine whether generalizations of Michaelis-Menten kinetics (\hyperref[intoEqtn:mm]{Equation \ref{intoEqtn:mm}}) \cite{Liebermeister:2006fm} are sufficient to quantitatively related physiological concentrations of substrates and enzymes to variable reaction flux.  To the extent that such an agreement is not possible, I then assessed whether metabolic regulation through allostery is able meaningful improve fit (\hyperref[introFig:primarytopics]{Figure \ref{introFig:primarytopics}B}).  This approach, called SIMMER (\underline{S}ystematic \underline{I}nference of \underline{M}eaningful \underline{M}etabolic \underline{E}nzyme \underline{R}egulation) is novel in that it is able to to scalably identify metabolic regulation. Accordingly, this approach was used to both confirm the physiological important of canonical yeast metabolic regulation and to identify novel regulation which was subsequently experimentally validated. By inferring reaction kinetics based on physiological variation in reaction species, SIMMER also naturally indicates the relative influence of substrates, products, enzymes and regulators in changing reaction flux. This suggests that metabolic activity is largely determined by changes in substrate and product concentrations based on local metabolic context, with direct regulation through regulated changes in protein abundance and allostery primarily acting to fine-tune metabolic behavior.

By deciding to directly measure protein abundance across the chemostat conditions, rather than relying upon transcript abundance as a surrogate for protein abundance I was guarding against possible deviations between protein and transcript abundance. Once protein relative abundance was measured, I was able to compare transcriptional changes to corresponding changes in protein abundance to assess whether the addition of proteomics was necessary (\hyperref[introFig:primarytopics]{Figure \ref{introFig:primarytopics}C}). The results of this analysis are presented in \hyperref[ch:pt_compare]{Chapter \ref{ch:pt_compare}}.  While strong transcriptional changes do indeed propagate into a corresponding change in protein abundance, additional regulation appears to be operating at the post-transcriptional level to further attune transcriptional responses to the metabolic environment. Regulation of protein degradation appears to be particularly important during nitrogen-limitation, where ribosomes are preferentially degraded and metabolic enzymes are selectively retained.

Code supporting each chapter is maintained on \href{https://github.com/shackett}{github}, access to private repositories is available upon request.

\begin{figure}[h!]
\floatbox[{\capbeside\thisfloatsetup{capbesideposition={right,top},capbesidewidth=4cm}}]{figure}[\FBwidth]
{\caption[Summary of primary thesis topics]{\fixedspaceword{\\Summary of primary thesis topics. \textbf{A)} A multi-omic data was generated involving four distinct classes of high-dimensional data. Previously measured transcript and metabolite relative abundances \cite{Brauer:2008jn, Boer:2010fb} were supplemented with measurements of enzyme relative abundances and inference of metabolism-wide fluxes. Each data type was measured across 25 standard chemostat conditions. \textbf{B)} Integration of metabolite, enzymes and flux measurements allows testing of reaction-level models of kinetics. Each possible model of kinetics is described as a non-linear function of substrates, products, enzymes, possible regulators and constants. Alternative kinetic models are compared based on model plausibility, $Pr(Model)$, and how well measured fluxes agree with the Michaelis-Menten kinetics, $Pr(Data | Model)$. \textbf{C)} The relationship between changes in transcript abundance and protein abundance can be modeled using differential equations. Departures between protein and transcript changes may reflect additional post-transcriptional regulated due to variable translational efficiency or protein degradation.
}}\label{introFig:primarytopics}}
{\includegraphics[width=\textwidth-4cm]{ch-intro/Figures/globalQuant.pdf}}
\end{figure}



\chapter{Related Work\label{ch:pastwork}}

Everyone needs a chapter about related work, so here is a placeholder.

% include other files for sections of this chapter. These use the 'input' command since each section within a chapter should not start a new page.
% If you want to swap the order of sections, it is as simple as reversing the order you include them. 
\section{Tables}
\label{sec:pastwork:tables}

Tables are also quite important. Any table that can fit entirely on one page can be a floating table. If a table is longer and will span multiple pages, a long table can be inserted in-line with the text. This is demonstrated in Table~\ref{tab:usage:options}, and explained in Appendix~\ref{ch:implementation}.

Tables that fit on one page use normal floating figures. Keep the 'p' placement option (in addition to 'h' and 't') so that if the float cannot fit in-line with the document text, it can be on a separate page by itself immediately after it is placed. Without the 'p' option, the float may get pushed to the end of the chapter, along will all other floats in the chapter that follow it.

Table~\ref{table:pastwork:publishing} lists the various options for publishing your dissertation, with costs, as of 2010. You will have to bring a check for the appropriate amount, made out to ``Princeton University Library'', when you submit your bound dissertation copies to Mudd Library, along with the appropriate forms and the electronic copy of your dissertation burned to a CD (not a DVD) as a single PDF file. (See~\cite{muddthesis2009}.)

Traditional publishing is cheaper initially and lets you earn royalties if the publisher sells many copies of your dissertation. However, most of us won't have a best-seller dissertation and most likely won't earn royalties anyway. Instead, by choosing open access publishing, your dissertation will be available online for free to anyone who is interested. I strongly advocate for open access, to maximize the impact of your research.

Your dissertation is protected by copyright regardless of whether or not you have the copyright registered. However, registration establishes a public record of your copyright claim~\cite{muddthesis2009}. ProQuest will submit the copyright registration for an extra fee (about \$55). Alternatively, you can register it yourself at the Copyright Office's website for only \$35: \url{http://www.copyright.gov/eco/}.

\begin{table}[htbp]
\centering
\caption[Thesis Publishing Options]{Thesis publishing options~\cite{mudd2009}, as of May 2010. }
\label{table:pastwork:publishing}
\begin{tabular}{p{0.3\textwidth} p{0.15\textwidth} p{0.15\textwidth} p{0.15\textwidth} p{0.15\textwidth}}
\toprule
\textbf{Publishing Method} & \textbf{Publishing Fee}
 & \textbf{Diploma Fee} & \textbf{Copyright Registration Fee} & \textbf{Total} \\
\midrule
\multicolumn{5}{c}{Traditional Publishing}\\
\midrule

Traditional without copyright registration
& 65 & 15 & -- & 80 \\[0.2em]

Traditional with copyright registration
& 65 & 15 & 55 & 135 \\[0.2em]

\midrule
\multicolumn{5}{c}{Open Access}\\
\midrule

Open access without copyright registration
& 160 & 15 & -- & 175 \\[0.2em]

Open access with copyright registration
& 160 & 15 & 55 & 230 \\

\bottomrule
\end{tabular}
\end{table}
\section{Figures}
\label{sec:pastwork:figures}

Everyone needs floating figures in their dissertation. 

As shown in Figure~\ref{fig:pastwork:titlepage}, the Mudd Library dissertation requirements~\cite{muddthesis2009} specify additional options for formatting the title page. For example, if your thesis has multiple volumes, or to indicate the proper formatting for a master's thesis.

\begin{figure}[htb]
  \begin{center}
    \includegraphics[width=0.9\linewidth]{ch-pastwork/figures/titlepage}
    \caption[Sample Title Page Layout]{Sample title page layout~\cite{muddthesis2009}}
    \label{fig:pastwork:titlepage}
  \end{center}
\end{figure}




\chapter{Usage\label{ch:usage}}

To start, in your main .tex file, use this class as your main documentclass instead of `report' or `book'. For example:
\begin{quote}
$\backslash documentclass[12pt,lot, lof]\{puthesis\}$
\end{quote}

In this example, we setup our document to use the PU Thesis style, with 12pt font for body text, and to include a List of Tables and List of Figures in the front matter. You could instead set an 11 point or 10 point font by changing the first option. You can also add `los' to include a list of symbols.

To use single spacing, add the option `singlespace'. This is a special option for the \texttt{puthesis} documentclass, which sets single spacing for both the front matter and for the document itself. Additional parameters should be set in your main .tex file, and are described in detail in Section~\ref{sec:usage:options}.

The template itself declares two other options, to be set immediately after the \texttt{documentclass} command. First is `printmode', declared with the command:
\begin{quote}
$\backslash newcommand\{\backslash printmode\}\{\}$
\end{quote}
This command, used later in the thesis.tex file, turns off the \texttt{hyperref} package and all internal links in the PDF file. This removes any colored links and highlighting that would not be appropriate in a printed and bound thesis. Instead the \texttt{url} package is loaded, so that \\url commands in your document will continue to work and urls will break properly across multiple lines.

When `printmode' is not specified, the hyperref package is included. It creates colored links for citations, footnotes, and internal references, which can be used to navigate the PDF document more easily. It also adds bookmarks to the PDF file, mirroring the table of contents. By default, it is set to use colored links. For the PDF file that you will submit electronically to ProQuest, this may not be desirable since some copies may be printed, while others will be used electronically. Thus another option, `proquestmode', is defined that keeps hyperref but disables colored links:
\begin{quote}
$\backslash newcommand\{\backslash proquestmode\}\{\}$
\end{quote}
This mode has no effect when used in combination with `printmode'. 


\section{Options}
\label{sec:usage:options}

In this section, we describe the options you can set when using this thesis class.
\tablespacing
% tablespacing is defined by the class to set single spacing for the long table when in doublespacing mode. If the singlespace option is set, this command has no effect.

\begin{longtable}{p{0.3\linewidth} p{0.6\linewidth}}

  % First page heading
  \caption[Options Provided by the PUthesis Class]{List of options for the puthesis document class and template} \label{tab:usage:options}\\
  \toprule
  \textbf{Option} & \textbf{Description} \\
  \midrule
  \endfirsthead

  % Future page heading
  \caption[]{(continued)}\\
  \toprule
  \textbf{Option} & \textbf{Description} \\
  \midrule
  \endhead

  % Page footer
  \midrule
  \multicolumn{2}{r}{(Continued on next page)}\\
  \endfoot

  % Last page footer
  \bottomrule
  \endlastfoot

  12pt &
  Specify the font size for body text as a parameter to \texttt{documentclass}. The Mudd Library requirements~\cite{muddthesis2009} state that 12pt is preferred for serif fonts (e.g., Times New Roman) and 10pt for sans-serif fonts (e.g., Arial).
  \\

  letterpaper &
  If your document is coming out in a4paper, your LaTeX defaults may be wrong. Set this option as a parameter to \texttt{documentclass} to have the correct 8.5"x11" paper size.
  \\

  lot &
  Set this option as a parameter to \texttt{documentclass} to insert a List of Tables after the Table of Contents.
  \\


  lof &
  Set this option as a parameter to \texttt{documentclass} to insert a List of Figures after the Table of Contents and the List of Figures.
  \\

  los &
  Set this option as a parameter to \texttt{documentclass} to insert a List of Symbols after the Table of Contents and the other lists.
  \\

  singlespace &
  Set this option as a parameter to \texttt{documentclass} to single space your document. Double spacing is the default otherwise, and is required for the electronic copy you submit to ProQuest. Single spacing is permitted for the printed and bound copies for Mudd Library.
  \\
  
  draft &
  Set this option as a parameter to \texttt{documentclass} to have \LaTeX mark sections of your document that have formatting errors (e.g., overfull hboxes). 
  \\

  % the cmidrule here spans both columns but is indented slightly on the left and right. 
  \cmidrule[0.1pt](l{0.5em}r{0.5em}){1-2}

  \raggedright
  $\backslash newcommand$ $\{\backslash printmode\}\{\}$ &
  Insert this command after the \texttt{documentclass} command to turn off the hyperref package to produce a PDF suitable for printing.
  \\

  \raggedright
  $\backslash newcommand$ $\{\backslash proquestmode\}\{\}$  &
  Insert this command after the \texttt{documentclass} command to turn off the `colorlinks' option to the hyperref package. Links in the pdf document will then be outlined in color instead of having the text itself be colored. This is more suitable when the PDF may be viewed online or printed by the reader.
  \\

  $\backslash makefrontmatter$ &
  Insert this command after the \texttt{$\backslash begin\{document\}$} command, but before including your chapters to insert the Table of Contents and other front matter.
  \\
  
  \cmidrule[0.1pt](l{0.5em}r{0.5em}){1-2}

  $\backslash title$ &
  Set the title of your dissertation. Used on the title page and in the PDF properties.
  \\

  $\backslash submitted$ &
  Set the submission date of your dissertation. Used on the title page. This should be the month and year when your degree will be conferred, generally only January, April, June, September, or November. Check the Mudd Library rules~\cite{mudd2009} for the appropriate deadlines.
  \\

  $\backslash copyrightyear$ &
  Set the submission year of your dissertation. Used on the copyright page.
  \\

  $\backslash author$ &
  Your full name. Used on the title page, copyright page, and the PDF properties. \\

  $\backslash adviser$ &
  Your adviser's full name. Used on the title page. \\

  $\backslash departmentprefix$ &
  The wording that precedes your department or program name. Used on the title page. The default is ``Department of'', since most people list their department and can leave this out (e.g., Department of Electrical Engineering), however if yours is a program, set $\backslash departmentprefix\{Program in\}$ \\

  $\backslash department$ &
  The name of your department or program. Used on the title page. \\

  \cmidrule[0.1pt](l{0.5em}r{0.5em}){1-2}
  
  \raggedright  
  $\backslash renewcommand$ $\{\backslash maketitlepage\}\{\}$ &
  Disable the insertion of the title page in the front matter. This is useful for early drafts of your dissertation. \\

  \raggedright  % full justification places the * in an awkward place
  $\backslash renewcommand*\{\backslash makecopyrightpage\}\{\}$ &
  Disable the insertion of the copyright page in the front matter. This is useful for early drafts of your dissertation. \\

  \raggedright 
  $\backslash renewcommand*\{\backslash makeabstract\}\{\}$ &
  Disable the insertion of the abstract in the front matter. This is useful for early drafts of your dissertation. \\

\end{longtable}
\bodyspacing
% bodyspacing restores double spacing or single spacing after the table

% need blank space after \bodyspacing

I've seen other people print their dissertations using $\backslash pagestyle\{headings\}$, which places running headings on the top of each page with the chapter number, chapter name, and page number. This documentclass is not currently compatible with this option -- the margins are setup to be correct with page numbers in the footer, placing them 3/4" from the edge of the paper, as required. If you wish to use headings, you will need to adjust the margins accordingly.
 





\chapter{Conclusion \& Future Work\label{ch:conclusion}}

Future Work

\section{Future Directions}

Future work should include options in the template for a masters thesis or an undergraduate senior thesis. It should also support running headings in the headers using the `headings' pagestyle.  The print mode and proquest mode included in the template might also be candidates to include in the class itself. 

\begin{outline}
\1 The use of genetic and environmental perturbations - genetic and environmental perturbations alter networks in distinct ways
\1 

\end{outline} % Conclusion
\section{Future Work}

The work that I have presented here can serve as a template for how metabolism and protein expression can be efficiently interrogated. My analysis of these questions, however, is limited by the absence of crucial information and an incomplete survey of physio-genetic variation. As such, when discussing future directions, I will primarily focus on how the caveats limiting my current research could be best addressed. From there, I will touch upon the real world utility of metabolic and expression models, specifically focusing on how such models could be applied to metabolic engineering to improve the efficiency of microbial compound production.

\subsection{Achieving a quantitative understanding of metabolism}

As discussed in \hyperref[ch:simmer]{Chapter \ref{ch:simmer}}, the kinetics of metabolic reactions can be analyzed \textit{en masse}. This is accomplished by measuring reaction flux and all reaction species across a sufficient array of conditions where alternative models can be statistically discriminated. The success of such an approach requires that one of the tested models approximates true \textit{in vivo} kinetics and that no measurements are systematically biased. In practice, no such guarantees are afforded at this juncture because a substantial list of caveats currently hinder SIMMER. These limitations largely arise from phenomenon with little demonstrated metabolic significance and unknown impacts. By experimentally addressing each question within the SIMMER framework, however, a more complete understanding of how metabolism is regulated can be found.

\subsubsection{Metabolomic incompleteness}

One current limitation of the SIMMER methodology is that while metabolomics is able to measure a meaningful fraction of important small molecules, this number is minute when compared to the total number of species in most metabolic reconstructions (e.g. 1,967 species reported in the yeast consensus reconstruction v7 \cite{BenjaminDHeavner:2013bp}). Mass spectrometry-based metabolomics is best at measuring abundant intracellular metabolites that readily ionize. This results in the preferential ascertainment of metabolites in central carbon metabolism, where, SIMMER can consequently be best applied. Because of the incompleteness of the ascertained metabolome, for most reactions, the primary substrate(s) have not been measured; and thus, SIMMER cannot be intelligently applied. In other cases, unmeasured products or minor substrates (such as protons or water) may prevent the identification of an appropriate kinetic model or lead to false positives. For example, a false positive may arise from the correlation between a missing product and a measured putative inhibitor. When such species are missing it is assumed that their concentration is invariant; this assumption, thus, will only impact the application of SIMMER to the extent that physiological variation in missing species impacts flux.

Measuring a wider array of metabolites would enable the broader application of SIMMER. A partial solution to this limitation will be naturally afforded through the passage of time; as mass spectrometer sensitivity continues to advance at a rapid pace \cite{Radionova:2015gs}, measurement of increasingly scarce species becomes more routine. Improvements in sensitivity are only a partial solution. Alternative protocols may be necessary to discriminate chemically similar species or to measure species that are not reliably detected using standard methods \cite{Lammerhofer:2013tz}.

Direct measurement of small molecules is the best solution to account for each small molecule's kinetic impact. Nevertheless, for some species measurement will be challenging due to their low concentration (e.g non-storage metabolites, outside of central carbon metabolism); lability (e.g. oxaloacetate); or measurement difficulty (e.g. protons and water). In some cases, such as for water, it might be reasonable to treat metabolites as invariant, but an interesting alternative solution to this limitation is the use of \textit{optimal regulators} (\hyperref[simmer_hyporeg]{Section \ref{simmer_hyporeg}}). Using this approach, both reaction kinetics and an optimal pattern of regulator abundance (informed by the metabolomics principle components) are simultaneously inferred. An analogous approach could be used to impute the abundance of an unmeasured metabolite, either for a single isolated reaction or for multiple coupled reactions.  For instance, the relative abundance of oxaloacetate could be inferred based on its shared involvement in mitochondrial malate dehydrogenase, citrate synthase and pyruvate carboxylase. Similarly, unknown proton concentrations could be assessed by aggregating numerous protonation reactions. 

While intriguing, jointly inferring the kinetics of coupled reactions detracts from the scalability of the SIMMER. A possible middle ground may be found by using an Expectation-Maximization-like approach. Such an approach would iterate between solving metabolite abundances using coupled reactions and then optimizing the kinetics of individual reactions conditioned on imputed metabolite abundance. On the more exploratory side, this PC-based imputation approach could be used to explore the kinetics of additional reactions. Multiple under-characterized reactions with unmeasured metabolites could be combined into a coupled inference problem. This could be a powerful approach for metabolic exploration; inference would still be underdetermined and the use of a bayesian approach would appropriately convey uncertainty in individual kinetic parameters and metabolite loadings. PC-based imputation may serve as a powerful gap-filling approach, not only in metabolism, but for elsewhere in genomics and beyond.

\subsubsection{Improving the precision of flux estimation}

In order to break metabolic pathways into reaction-level inference problems, reaction flux must be treated as a fixed quantity. Determining flux is challenging; fluxes are rates rather than physical species that can be directly measured. Techniques to determine these rates have either largely estimated fluxes that maximize growth using Flux Balance Analysis (FBA) \cite{Orth:2010hb} or introduced isotopically labelled metabolites that propagate through metabolism using Metabolic Flux Analysis (MFA) \cite{Zamboni:2009bp}. The needs of inferring fluxes in chemostat culture are somewhat different from standard problems. Cells grow at known rates and metabolic sources and sinks can be readily estimated. Using this information, the variant of Flux Balance Analysis (FBA) that I developed is well-suited for estimating fluxes that are strongly determined by these boundary fluxes. This method focuses on finding a least-squares solution that optimally conforms to experimental data by using quadratic programming, 

It is unlikely that all fluxes are well-estimated by virtue of metabolite uptake, excretion and incorporation into biomass components. Based on boundary fluxes alone, it is difficult to estimate flux through pathways that have stoichiometrically equivalent routes for flux (such as in 1C metabolism) or through pathways that primarily affect energy or reducing equivalents rather than uptake, excretion or biomass component synthesis. To better quantify such pathways, future work should supplement experimental boundary fluxes. These can be determined through spectrophotometry and $^{1}$H-NMR with additional experimental measurements to further constrain possible flux.

Respirometry could allow for better estimation of glycolytic and TCA flux is respirometry.  In microbes, oxygen consumption is a surrogate for electron transport chain flux. The respiratory quotient $\left(\sfrac{\nu_{CO_{2}}}{\nu_{O_{2}}}\right)$ is a strong indication of the balance of fermentation versus aerobic respiration \cite{Boer:2003fi,BARFORD:1979ei}. To further refine estimates of intracellular flux, techniques from MFA could be applied to chemostat culture. This could be particularly useful for determining the ratio of fluxes between pathways, such as between glycolysis and the oxidative pentose phosphate pathway, or between lower-glycolysis and the TCA cycle \cite{Sauer:2006ii, Jazmin:2013fg, Kromer:2014wr}. MFA is generally considered a comprehensive framework for the estimation of metabolism-wide fluxes that can incorporate experimental boundary fluxes. But, estimates of reaction flux could also be incorporated into a quadratic programming based method.

Estimating additional reaction fluxes using respirometry \&/or MFA would allow for better estimation of the expected values of pathway fluxes; and, the uncertainty in these measurements could also be used to better characterize the uncertainty in metabolism-wide flux. The quadratic penalty is constructed by weighting least-squared deviations by experimental precision analogously to gaussian log-likelihood $\left(\nu^{T}\Sigma^{-1}\nu = \sum_{j = 1}^{J}\frac{\left(x_{j} - \mu_{j}\right)^2}{\sigma^{2}_j}\right)$. Accordingly, uncertainty in flux can be assessed by the maximum and minimum flux that could go through each reaction at the solution minimum as well as at sub-optimal solutions that approximate cuts through the likelihood distribution (\hyperref[conclusion-fluxUncertainty]{Figure \ref{conclusion-fluxUncertainty}}). This is conceptually similar to estimating confidence intervals based on the likelihood ratio. SIMMER currently deals with uncertainty in flux using flux variability analysis \cite{Mahadevan:2003wq}, based on how constrained individual fluxes are at the solution optimum. A solution that more fully addresses uncertainty could be readily incorporated into SIMMER by integrating over the full flux probability density function.  This solution was not used in \hyperref[ch:simmer]{Chapter \ref{ch:simmer}} because the absence of internal fluxes prevented strong constraints on most fluxes at sub-optimal values. This resulted in pathological flux distributions. 

\begin{figure}[h!]
\begin{center}
\includegraphics[width=0.8\textwidth]{ch-conclusion/fluxVar.pdf}
\caption[Procedure for accounting for flux uncertainty]{Procedure for accounting for flux uncertainty. \textbf{A)} The relationship between the maximum and minimum flux that are possible at a given quadratic penalty / log-likelihood can be used to construct a \textit{de facto} log-likelihood distribution for each reaction's flux. \textbf{B)} The log-likelihood distribution corresponds to a probability distribution, p(v).}
\label{conclusion-fluxUncertainty}
\end{center}
\end{figure}



\subsubsection{Mischaracterization of enzyme activity and kinetic differences between isoenzymes}

When incorporating enzyme-level information into SIMMER, measurement of protein relative abundance can be accurately applied in a simple case: (1) a single gene produces an enzyme, (2) each copy of this enzyme is kinetically indistinguishable. Outside of this simple scenario, kinetic inference is complicated by multiple distinct pools of enzymes that catalyze a single reaction. Relevant enzymatic pools with greatly varying kinetics include: isoenzymes that catalyze the same reaction \cite{Wilson:2003cb}, enzyme complexes where pools of active enzyme may not be reflected by the abundances of monomers \cite{Cohen:2000bj}, and post-translationally modified enzymes with altered metabolite affinity or activity \cite{Cohen:2000bj, Schulz:2014eo}. 

Because many enzymes do not operate as monomers, the abundance of each catalytically-active complex that catalyzes a given reaction should be included as a separate distinct ``enzyme'' in SIMMER. We have an incomplete understanding of how complex abundance and activity relate to the abundance of individual components. In general, multimers are stable, while a stoichiometric misbalance between components results in degradation of the excess \cite{Marianayagam:2004ie}. In such simple examples, the abundance of a complex would reflect the abundance of each component monomer.  The genome scale model used contains annotated complex membership \cite{BenjaminDHeavner:2013bp}, so to date, complex abundance has been estimated using the precision-weighted mean over all complex components. In some cases such a simplified view of complex formation may be inaccurate.  More sophisticated models of complex formation that estimate complex abundance from proteomics data could easily be incorporated into SIMMER, however.
 
When multiple complexes or isoenzymes catalyze a single reaction, we are left with two options; we can either assume that each set of enzymes is kinetically similar or different. If we assume that all pools of enzymes interact with metabolites in a similar manner (shared $k_{d}$ and possible $k_{i}$/$k_{a}$), the occupancy of each enzyme will be the same ($\sfrac{\nu}{\nu_{max}}$). $k_{cat}$ will vary to reflect that measured enzyme concentrations are relative and sizes of enzyme pools may differ. This is the standard approach currently used in SIMMER. This minimizes the total number of fitted kinetic parameters; and in most cases, a single complex likely carries the bulk of flux. If we chose to relax this assumption, SIMMER is already  compatible with reaction forms in which one or more kinetic parameters (besides $k_{eq}$) differ between isoenzymes or complexes. Treating simple Michaelis-Menten kinetics catalyzed by two isoenzyme (E$_{a}$ and E$_{b}$) as an example, reaction flux would follow \hyperref[Eq:isoenzymeMM]{Equation \ref{Eq:isoenzymeMM}}. DAHP synthase (Aro3/Aro4), the first step of aromatic amino synthesis, is one clear example in yeast where such isoenzyme-specific regulation is likely relevant; Aro3p is specifically inhibited by phenylalanine, while Aro4p is specifically inhibited by tyrosine \cite{Schnappauf:1998ec}. The kinetics of this reaction tested using isoenzyme-specific regulation, but the support for this regulation was quantitatively weaker than predicted regulation by phenylpyruvate (a phenylalanine precursor) (\hyperref[signif_regulators]{Table \ref{signif_regulators}}). 

\begin{equation}
\nu = k_{cat}^{E_{a}}\left[E_{a}\right]\frac{\left[S\right]}{\left[S\right] + k_{m}^{E_{a}}} + k_{cat}^{E_{b}}\left[E_{b}\right]\frac{\left[S\right]}{\left[S\right] + k_{m}^{E_{b}}}\label{Eq:isoenzymeMM}
\end{equation}

Pools of enzymes may differ kinetically because they are composed of different proteins, but the kinetics of otherwise identical enzymes or complexes may also differ if they possess different post-translational modifications.  Post-translational modifications (particularly phosphorylation) are rampant in yeast \cite{Fiedler:2009hx}; their effects on metabolism, however, are largely unknown aside from a few phosphorylation events with an implicated kinetic role, particularly in the TCA cycle \cite{Schulz:2014eo}. While post-translational modifications have not been characterized under the measured conditions, SIMMER could be easily adapted to suggest the physiological role of post-translational modifications in a similar manner to testing allostery. To illustrate how such inference could be conducted, we can consider whether assessing the effect of a single phosphorylation event with a putative effect on enzyme activity explains measured flux significantly better than a model without regulation. If we knew the relative abundance of both the modified protein, $\left[E_{p}\right]$, and unmodified protein $\left[E_{u}\right]$, reaction kinetics would be described by \hyperref[Eq:ptmMM1]{Equation \ref{Eq:ptmMM1}}. This equation fits one more parameter than the minimal model.

\begin{equation}
\nu = k_{cat}^{p}\left[E_{p}\right]\frac{\left[S\right]}{\left[S\right] + k_{m}} + k_{cat}^{u}\left[E_{u}\right]\frac{\left[S\right]}{\left[S\right] + k_{m}}\label{Eq:ptmMM1}
\end{equation}

In general, we will not have measured both the modified and unmodified protein. Rather, we will have measured the abundance of a covalently-modified peptide, $\left[E^{p}_{c}\right]$, and the relative abundance of the whole protein ($\left[E^{p}_{t}\right]$: unmodified plus modified). With this input data, we cannot distinguish scenarios where the maximum fraction of covalent modified peptides is 2\% of the total protein, from when 50\% of the protein is modified. To deal with this unknown, the maximum labelling fraction, $\phi_{max} \in [0, 1]$ must be estimated concurrently with kinetic parameters that post-translational modifications are posited to effect (\hyperref[Eq:ptmMM2]{Equation \ref{Eq:ptmMM2}}).

\begin{align}
c^{*} &= \underset{c\in C}{argmax} \sfrac{\left[E^{p}_{c}\right]}{\left[E^{t}_{c}\right]}\notag\\
\phi_{c} &= \phi_{max}\left(\frac{\sfrac{\left[E^{p}_{c}\right]}{\left[E^{t}_{c}\right]}}{\sfrac{\left[E^{p}_{c^*}\right]}{\left[E^{t}_{c^*}\right]}}\right)\notag\\
\nu_c &= k_{cat}^{p}\left[E_c^{p}\right]\frac{\left[S_c\right]}{\left[S_c\right] + k_{m}} + k_{cat}^{t}(1-\phi_c)E_{c}^{t}\frac{\left[S_c\right]}{\left[S_c\right] + k_{m}}\notag\\
 &= \left(k_{cat}^{p}\left[E_c^{p}\right] + k_{cat}^{t}(1-\phi_c)E_{c}^{t}\right) \frac{\left[S_c\right]}{\left[S_c\right] + k_{m}}\label{Eq:ptmMM2}
\end{align}

\subsubsection{Is the ``well-mixed'' assumption appropriate?}

Biochemistry is greatly simplified when we assume that processes are spatially uniform and thus ``well-mixed.'' Assuming that all cells within a culture are effectively equivalent is often an obligate requirement to collect enough physical material to quantify biological species. The use of many cells decreases the noise inherent in individual cells but removes information about how cells meaningfully differ. 

In extreme cases, culture-wide behavior may not reflect the behavior of individual cells. As an example, consider a metabolic reaction: A + B $\rightarrow$ C. If the metabolite A is present in only half of the cells, and B is present in the complementary half, C will not be produced in either cell population. However, applying the culture-wide well-mixed assumption interaction between A and B would be possible, allowing the pathological synthesis of C to proceed. An absence of \textit{in vivo} colocalization can lead to the same fallacy when either cellular compartmentalization is ignored or enzymes \&/or metabolites are erroneously assumed to be evenly distributed within a compartment.

In some cases, the assumption of cell-level similarity is likely appropriate, at least to a first approximation. The best-case scenario for such an assumption to hold is likely during chemostat growth, during which many generations of similar environmental conditions minimize the impact of environmental heterogeneity and historical contingency. The culture-wide reproducibility of chemostats suggests that the biological variance between individual cells is minimal. Although it is routine to verify the population is morphologically similar, at least in terms of cellular volume, during chemostat culture. While chemostat growth may be a best-case scenario for minimizing cell-to-cell variation, both random \cite{BarEven:2006dz, Kaern:2005gr} and non-random variation at the single cell level may be important at the population-level.

The clearest example of a process which varies between cells is the temporal evolution of the mitotic cell-cycle whose sinusoidal oscillations could not be summarized by a bulk-sum analysis \cite{Hartwell:1974uy, Spellman:1998wj}. Such temporally variable processes are not readily observed unless cellular population are coordinated\cite{Hartwell:1974uy, Tu:2006cl} or single cells are tracked through morphology \cite{Herskowitz:1988ut}, microscopy \cite{Venturelli:2015ec} or single-cell sequencing \cite{Patel:2014dt}. Each of these latter techniques has revealed that otherwise similar cellular populations are frequently composed of distinct clusters of cells. Nonetheless, the scope of such heterogeneity remains unclear. If single-cell concentrations of reaction species were available, SIMMER could be applied to model flux heterogeneity, but gaining the necessary input data is not currently possible. While advances in microscopy and automation enable the simultaneous tracking of larger numbers of proteins \cite{Ghaemmaghami:2003ds, Dubuis:2013cw}, understanding of single cell metabolic heterogeneity is currently limited by inadequate methods to detect metabolite levels of single cells \cite{Zenobi:2013il}.

As previously discussed, it is routine to measure the culture-wide abundance of biological species without accounting for intercellular variability. In general, such measurements are also carried out without regard for intracellular heterogeneity, either at the scale of compartments or even finer-scale localization. Eukaryotic cells are composed of distinct compartments with their own complement of metabolites and proteins, which jointly govern the compartment's metabolic activity. Within each compartment, metabolic activity may be further impacted by colocalization of enzymes, allowing for efficient metabolic channeling of metabolites between sequential enzymatic steps \cite{Ovadi:1995wy}. To appropriately account for intracellular heterogeneity, metabolic species should be measured at their \textit{in vivo} concentration. Perhaps intracompartmental variable concentrations could be systematically incorporated into genome-scale models through appropriate microscopy experiment. Such concerns are secondary, however, to accounting for compartment-wise metabolic variation.

The absence of compartment-scale resolution to metabolic data does not necessary greatly impact SIMMER. Let us consider the concentration of a metabolite that is only found in a single compartment. If the compartments' volume is proportional to intracellular volume then the metabolite's intracellular concentration is proportional to an absolute abundance measurement. This results in an inferred kinetic constant off by a multiplicative constant, although the \textit{in vivo} reaction kinetics can still be determined. If, however, metabolites are differentially partitioned across conditions (between the cytosol, mitochondria and vacuole in particular), relative concentrations in compartments will not reflect that of whole-cells \cite{Kitamoto:1988wc}.

Differential partitioning could explain why phenylpyruvate was erroneously predicted as an inhibitor of DAHP synthase rather than the true regulator phenylalanine \cite{Schnappauf:1998ec}. Phenylalanine is partially stored in the vacuole, while phenylpyruvate is likely primarily cytosolic \cite{Kitamoto:1988wc}. Because phenylalanine is made from phenylpyruvate in the cytosol, phenylpyurvate may be a better proxy for cytosolic phenylalanine concentrations than whole-cell measurements. In addition to the vacuole, estimating the concentrations of mitochondrial metabolites is also challenging. This is because relative mitochondrial volume may vary across conditions, as suggested by the strong covariation of mitochondrial enzymes (Supplemental Figure 2) and by the fact that many mitochondrial metabolites are also present in the cytosol but do not readily mix.  

To better address the problems apparent when compartmentalization is ignored, in eukaryotes, each major metabolic compartments should be treated with the same care as with separate experimental conditions. Techniques are emerging that allow for analysis of both metabolomics and proteomics of individual sub-cellular compartments \cite{Klie:2011kq, Wuhr:2014fr}. Because these methods are relatively gentle, it is unclear whether such approaches could be used to quantify the mitochondrial metabolome due to its great level of activity. Such methods may be necessary to kinetically model the yeast TCA cycle. A clearer investigation of compartment-specific biomolecules should be complemented by microcopy-based approaches for measuring compartment volumes \cite{JENSEN:1993bz, Ghaemmaghami:2003ds}.

\subsubsection{Evaluating modeling assumptions}

I have previously discussed the numerous ways in which culture-wide measurements of reaction species and modeling assumptions may limit our ability to accurately determine reaction kinetics. These limitations are primarily experimental rather than conceptual. Measurements of post-translational modification and compartment-specific quantification of metabolite abundances could be readily incorporated into SIMMER. Insufficient data greatly limits the number of reactions whose kinetics can be interrogated. For other reactions, our inability to describe reaction kinetics despite measuring the primary reaction species  and possible regulators suggests that other factors are kinetically important. Missing data could limit our ability to identify appropriate reaction kinetics or could even lead to false-positive predictions of regulation. Nevertheless, our ability to accurately predict reaction flux through numerous reactions using only reaction species and regulators suggests that our simplified models are more likely appropriate than not.

From my analysis, for 44 of 56 tested reactions, simple Michaelis-Menten kinetics including at most two regulators can accurately relate metabolite and enzyme concentrations to reaction flux. This is impressive firstly due to the large number of reasons why this could fail. Secondly, mechanistic modeling of the kinetics of well-studied reactions rarely reduces to Michaelis-Menten kinetics, and often many regulators are thought to collectively regulate activity \cite{Hill:1977vm}.  Under assay conditions, significant deviations from Michaelis-Menten kinetics may exist and the activity of regulators can be demonstrated. For the purpose of understanding flux control, a reaction form does not need to be principled mechanistically; rather, it must only account for an approximately correct relationship between species and resulting flux across the physiological conditions investigated \cite{Fell:1997wg}. 

When trying to estimate the kinetics of reactions that cannot be described using metabolites, enzymes and flux, additional experimental data will surely shed light on the role of unaccounted for sources of regulation. In other cases, however, inferring appropriate reaction kinetics may require alternative modeling assumptions rather than additional data. Such modeling assumptions could be easily relieved; in principle, using SIMMER any relationship between reaction species and resultant flux can be evaluated using flux. If an alternative reaction mechanism could be significantly supported over alternative contenders, this would suggest validity analogously to how reaction mechanisms are tested \textit{in vitro} \cite{CornishBowden:2012wb}. Other kinetic phenomenon such as cooperativity \cite{Hill:1910vo, Bush:2012jf} can also be tested, as was demonstrated when evaluating allosteric cooperativity.

\subsubsection{Application of SIMMER using more comprehensive data, broader conditions and new organisms}

Because we have characterized a reproducible set of conditions that contains much of meaningful metabolic variation in yeast, I hope this dataset will serve as an important scaffold for future work. In order to better understand yeast metabolism, it will be important to refine and expand our input data in the ways that I have discussed above. It will also be crucial to expand the number of experimental conditions used to study kinetics. Investigating additional conditions is important for two reasons. First, when too few conditions have been tested, alternative models of reaction kinetics may not be discriminated, for instance, due to correlation of putative regulators. Second, we can only identify important regulation that impacts the conditions that we have studied. By expanding the number of conditions that we investigate, we can unmask regulation that is important in these conditions. Furthermore, SIMMER can identify this regulation by virtue of the kinetic importance it gains in these new conditions.

To improve kinetic inference when choosing additional conditions to study, we would ideally maximize the number of distinct input-output pairs that can found (i.e. combinations of flux and all potential reaction species). We would also prioritize conditions spanning meaningful metabolic transitions. These two goals are broadly compatible; diverse conditions generate extensive multi-`omic variation and reveal metabolic regulation necessary to operate at these states.

For example, if we grew yeast using ethanol as a carbon source rather than glucose, this would result in a unique pattern of flux, enzymes and metabolites in addition to the emergence of conditionally-meaningful regulation. This would allow us to confirm the regulation that facilitates the glycolytic to gluconeogenic switch \cite{Zampar:2013fr}. It also would inform regulation more broadly through the addition of new metabolic state.

Similar logic could be used to better understand the specific metabolic impact of regulators of protein abundance. These regulators could be important transcription factors such as Gcn4 or affect signaling pathways such as Ras.  For example, a hypoactive variant of the Ras-regulator Ira2 in yeast greatly impacts transcript and protein abundance, but also alters the levels of glycolytic metabolites and flux \cite{Breunig:2014bu}. Through SIMMER's connection of $\left[regulator \rightarrow enzymes \rightarrow metabolites \rightarrow flux\right]$, the specific manner in which hypoactive Ira2 impacts flux could be discerned from ancillary changes in enzyme and metabolite. 

While SIMMER is a powerful way to understand yeast metabolism, its application to a well-studied, relatively simple organism does not fully highlight SIMMER's utility. Yeast metabolic regulation has been studied through decades of \textit{in vitro} biochemistry and genetics; however, such time-intensive characterization is not feasible for the majority of organisms. While previously reported \textit{in vitro} regulators can be prioritized and physiologically validated using SIMMER, in general, such prior assumptions are unnecessary or can be supplanted by regulation favored through genetic conservation.

When applying SIMMER to new organisms, or reactions with little or no previous chemical interrogation, the use of optimal regulators is a powerful way to prioritize regulatory metabolites. The reductionist approach of SIMMER will, moreover, be greatly useful when investigating the metabolic regulation of more complicated organisms, including humans. In higher eukaryotes, flux inference based on steady-state assumptions and parsimony largely break down; these methods are primarily supplanted by pathway-specific methods based on the dynamic turnover of isotopically-labelled metabolites. The scalability of SIMMER could help to elucidate the regulation and kinetics of these pathways. Because measurement of flux are routinely complimented with metabolomics and the generation of mutants that perturb pathway activity and metabolites, SIMMER could be applied largely with existing data.

\subsubsection{Translating reaction-level kinetics into metabolism-level predictions}

Determining reaction kinetics is important because it tells us how each reaction in a complex dynamical system will behave as function of its inputs. This is of massive importance in metabolic engineering. To increase flux through a pathway of interest that produces a commodity chemical, ideally we would know how pathway flux responds to changes in pathway enzymes. Due to the difficulty of reaction form inference, current models of metabolic regulation are limited by an incomplete understanding of metabolic regulation and reaction kinetics. Such models may lead to predicted metabolite-enzyme-flux relationships, which agree with the data on which they are fitted. They rarely, however, generalize when applied to new conditions. Because of these inaccuracies, it is difficult to predict how pathway flux will respond to genetic perturbations and even more challenging to understand how the metabolic network (and, indeed, broader regulation) will respond to such a perturbation. 

Due to these limitations, progress in metabolic engineering has generally been born out of brute force trial and error. Understanding how metabolism responds to changes in enzyme abundance is well-studied in Metabolic Control Analysis (MCA). Unfortunately, MCA generally requires a complete kinetic description of all relevant reactions; so, to date, its application has been largely impractical \cite{Kacser:1973fe, Fell:1997wg}. By facilitating the estimation of reaction forms, SIMMER is an ideal tool to inform MCA-based pathway models. To demonstrate the application of SIMMER to MCA, I generated an MCA-based model of glycolysis, which reproduces two primary aspects of yeast glycolytic control.

Similar methods can be used to predict changes in flux in response to genetic perturbations or the removal of allostery. This would entail introducing an instantaneous change in a parameter (enzyme abundance or removing allostery) and then following the time-evolution of the system as it reaches a new steady-state. Each reactions form is fitted across a wide range of metabolic conditions and contains relevant regulation. Accordingly, the behavior of individual steps should accurately respond to moderate changes, resulting in reasonably accurate predictions of pathway flux.

\subsection{Towards a comprehensive model of protein expression}

By analyzing systematic deviations between proteins and transcripts in \hyperref[ch:pt_compare]{Chapter \ref{ch:pt_compare}}, I suggested that post-transcriptional regulation provides a highly relevant layer of protein-expression control. This analysis was not ideal due to differences in experimental design between the two compared datasets. By measuring proteins and transcripts in the same culture, and ideally accounting for biological variance through replication, a more faithful comparison between these species could be conducted.  Such an analysis could also disambiguate the role of translational efficiency and protein degradation through experimental determination of one (or both) of these processes using existing methods \cite{Ingolia:2009dp, Belle:2006hv}.

The central dogma of molecular biology indicates that information in a cell largely flows from DNA to RNA to protein. The nearly universal relationship is limited because it is non-quantitative. Meanwhile, the composition of the proteome is governed by the variable rate of information transfer both between-genes and across-conditions. Towards this end, a quantitative central dogma of molecular biology is needed to account for the relationship between changes in protein expression and the intra- and extra-cellular environment.

Such a quantitative model of protein expression could be directly tied into models that depend on protein activity, such as metabolism \cite{OBrien:2013fl}.  One essential challenge when trying to predict the behavior of metabolism is that an induced change in protein expression will result in a change in flux and metabolite abundance. This will, consequently result in a secondary change in gene expression. This change could be due to the cost of over-expression \cite{Dykhuizen:1987uq}, impacts of nutrient sensing on transcriptional \cite{Brauer:2008jn} and post-transcriptional control (\hyperref[ch:pt_compare]{Chapter \ref{ch:pt_compare}}), or effects such as osmolarity and pH \cite{Csonka:1991wf}. When there are unknown responses to a genetic or environmental perturbation the accuracy of metabolic prediction declines as metabolism is pushed beyond measured conditions.

If an integrated expression-metabolism model were possible, metabolic behavior could be predicted as a function of the environment and expression. Expression could, in turn, be predicted based on environmental and metabolic state. The industrial impacts of this development would be massive. By reliably connecting genetic and environmental changes to their resulting metabolic impact, pathway fluxes could be efficiently altered; such fluxes could produce either native compounds or favor synthetic pathways that synthesize foreign compounds. In this case, pathway flux would operate as a tunable knob. Balancing flux into a desired pathway with fluxes that fulfill an organism's biosynthetic requirements would allow for predictable steady-state synthesis. 

Predicting metabolic behavior based on an organism's genetic and environmental state is the key to unlocking the industrial capabilities of microbes. Establishing this link between genotype and phenotype through context-dependent mechanistic models, is the same essential connection that is sought when investigating the genetic basis of complex traits and disease. Expanding this connection to higher-level traits could be accomplished using the same principles as those used to dissect the complexity of yeast. When studying disease, models of cellular behavior must first be improved. Next, models relating organ-level behavior must be developed. Ultimately, these systems must be integrated under that auspices of medicine, itself an endeavor which has worked to generate \textit{ad hoc} models of disease occurrence. 
 % Future work

\appendix % all chapters following will be labeled as appendices
\chapter{Implementation Details\label{ch:implementation}}

Appendices are just chapters, included after the $\backslash appendix$ command.

\section{Switching Formats}
When switching \texttt{printmode} on and off (see Section~\ref{sec:usage:options}), you may need to delete the output .aux files to get the document code to compile correctly. This is because the hyperref package is switched off for \texttt{printmode}, but this package inserts extra tags into the contents lines in the auxiliary files for PDF links, and these can cause errors when the package is not used.

\section{Long Tables}

Long tables span multiple pages. By default they are treated like body text, but we want them to be single spaced all the time. The class therefore defines a new command, $\backslash tablespacing$, that is placed before a long table to switch to single spacing when the rest of the document is in double spacing mode. Another command, $\backslash bodyspacing$, is placed after the long table to switch back to double spacing. Normal tables using \texttt{tabular} automatically use single spacing and do not require the extra commands.

When the documentclass is defined with the `singlespace' option, these commands are automatically adjusted to stay in single spacing after the long table.

Make sure there is always at least one blank line after the $\backslash bodyspacing$ command before the end of the file.

Some times long tables do not format correctly on the first pass. If the column widths are wrong, try running the \LaTeX compiler one or two extra times to allow it to better calculate the column widths.

If you want your long table to break pages at a specific point, you can insert the command $\backslash pagebreak[4]$, to tell \LaTeX that it really should put a page break there. $\backslash pagebreak[2]$ gives it a hint that this is a good place for a page break, if needed. If there's a row that really should not be broken across a page, use $\backslash \backslash *$, which will usually prevent a pagebreak. 

\section{Booktabs}
The booktabs package is included to print nicer tables. See the package documentation~\cite{fear2005booktabs} for more details and motivation. Generally, all vertical lines are removed from the tables for a better visual appearance (so don't put them in), and better spacing and line thicknesses are used for the horizontal rules. The rules are defined as $\backslash toprule$ at the top of the table, $\backslash midrule$ in between the heading and the body of the table (or between sections of the table), and $\backslash bottomrule$ at the end of the table. $\backslash cmidrule$ can be used with the appropriate options to have a rule that spans only certain columns of the table.

\section{Bibliography and Footnotes}

The bibliography and any footnotes can also be single spaced, even for the electronic copy. The template is already setup to do this.

Bibliography entries go in the .bib file. As usual, be sure to compile the \LaTeX code, then run BibTeX, and then run \LaTeX again.

To cite websites and other electronically accessed materials, you can use the `@electronic' type of BibTeX entry, and use the `howpublished' field to include the URL of the source material.

The formatting of bibliography entries will be done automatically. Usually the titles are changed to have only the first word capitalized. If you'd prefer to have your original formatting preserved, place the title in an extra set of curly braces, i.e., ``title = \{\{My title has an AcroNyM that should stay unchanged\}\},''.

\section{Figures and Tables}
The captions of figures and tables take an optional parameter in square brackets, specifying the caption text to be used in the Table of Contents. The regular caption in curly braces is used for the table itself.

Generally captions for tables are placed above the table, while captions for figures are placed below the figure.




\chapter{Printing and Binding\label{ch:printing}}

\section{Printing}

For the library copies of your dissertation, you must use archival quality printing and binding. This means acid-free paper, containing at least 25\% cotton fiber. Triangle Repocenter on Nassau Street in Princeton offers both 25\% cotton paper and 100\% cotton paper. Most people choose the 25\% cotton paper, and this is generally recommended by the binders. The 100\% copy paper is somewhat thicker and the extra expense is unnecessary. 

Triangle offers online submission of your printing and binding order at: \url{http://triangleprinceton.com/collegiatebinding/thesis/}. If you request binding from them, they will deliver the paper copies to Smith-Shattuck Bookbinding for you and allow you to pick up the completed copies at their store on Nassau Street. The whole process takes 2-3 business days, but check with them in advance during the busy thesis-printing season in April and May. 

Currently, your printed and bound dissertation copies can be single spaced. Only the electronic copy submitted to ProQuest must be double spaced. All copies must be printed single-sided, with specific margins. 

\section{Binding}

An archival-quality sewn binding is required for the library copies of your dissertation. Smith-Shattuck Bookbinding is highly recommended, and is used by most students. Triangle Repocenter will send your copies there for you, greatly simplifying the process, but you can call Smith-Shattuck with special requests. 

The ``library standard'' sewn binding is sufficient for the copies to be sent to Mudd Library. It uses a black buckram cloth cover, which is the most popular option. For extra copies for yourself and your family members, you can choose ``buckram roundback binding'', which adds decorative lines on the spine, and printing of the title and author on the front cover. For a small additional fee, you can include the Princeton University shield on the front cover and a ribbon bookmark. Leather covers are also available. See Smith-Shattuck's website for more details at: \url{http://www.thesisbookbinding.com/}. 


% Make the bibliography single spaced
\singlespacing
\bibliographystyle{plain}

% add the Bibliography to the Table of Contents
\cleardoublepage
\ifdefined\phantomsection
  \phantomsection  % makes hyperref recognize this section properly for pdf link
\else
\fi
\addcontentsline{toc}{chapter}{Bibliography}

% include your .bib file
\bibliography{thesis}

\end{document}

